\chapter{Ser Ninguém}

\lettrine{T}{entemos observar a cessação} ou o fim das coisas em pequena escala,
prestando atenção especial à parte final da expiração. Deste modo, no dia a dia,
percebemos finais comuns aos quais ninguém raramente presta atenção. Achei esta
prática muito útil porque é uma forma de se perceber a natureza mutável do reino
condicionado à medida que vivemos o dia a dia. Da maneira como vejo, foi para
esses estados mentais comuns que o Buddha estava a apontar, não para os estados
concentrados especiais bem desenvolvidos.

No primeiro ano que pratiquei, estava sozinho e consegui entrar em estados
mentais concentrados bem desenvolvidos, que realmente apreciei. Depois fui para
Wat Pah Pong onde a ênfase estava no modo de vida de acordo com a disciplina e
rotina do Vinaya. Lá, era sempre necessário sair para fazer a ronda da esmola
todas as manhãs e os cânticos de manhã e à noite. Quando se é jovem e saudável,
é suposto participar dessas rondas de esmolas muito longas -- havia rondas mais
curtas que monges mais idosos e fracos podiam fazer. Naquela época eu era muito
vigoroso, então fazia as longas, longas rondas de esmolas e depois voltava
cansado, em que de seguida dava"-se a refeição e à tarde todos tínhamos tarefas
para fazer. Não era possível, nessas condições, permanecer num estado
concentrado. A maior parte do dia era ocupado pela rotina da vida diária.

Então fartei"-me disso tudo, fui ter com Luan Por Chah e disse"-lhe: `Não
consigo meditar aqui', e ele começou a rir"-se de mim e a dizer para todos: `o
Sumedho não consegue meditar aqui!' Para mim a meditação era aquela experiência
extremamente especial que tinha vivido, adorado. Mas de imediato Luang Por Chah
estava obviamente a apontar para a normalidade do dia a dia, o levantar, as
rondas de esmolas, a rotina de trabalho, as tarefas domésticas. Tudo contribuía
para a consciência espiritual. E ele não parecia minimamente preocupado em
apoiar"-me nos meus desejos para poder ter uma forte experiência de privação
sensorial ao não ter que fazer todas essas pequenas tarefas diárias. Ele não
parecia importar"-se com isso; então acabei por me conformar e aprender a
meditar na rotina do dia a dia. E a longo prazo, isso tem sido o mais útil.

Nem sempre tem sido como eu quis, porque normalmente uma pessoa quer aquilo que
é especial: ter compreensões claras e maravilhosas, uma luz resplandecente,
felicidade, êxtase e arrebatamentos incríveis -- não só estar calmo e feliz --
mas no céu!

No entanto, reflectindo sobre a vida nesta forma humana: as coisas são
simplesmente como são, o podermo"-nos sentar e levantar tranquilamente e
contentarmo"-nos com o que temos; é isso que faz da nossa vida uma experiência
diária alegre e não dolorosa. E é assim que a maior parte da nossa vida pode ser
vivida -- não podemos viver em estados de êxtase e felicidade e lavar a louça ao
mesmo tempo, certo? Eu costumava ler sobre a vida de santos que estavam de tal
ordem embrenhados em êxtases, que não eram capazes de fazer nada a um nível
prático. Embora o sangue fluísse de suas palmas das mãos e eles conseguissem
realizar feitos que os fiéis apressariam"-se para ver, quando se tratava de algo
prático ou realista, eles eram bastante incapazes.

E ainda assim, quando contemplamos a disciplina do Vinaya em si, é um treino
para sermos conscientes. Trata"-se de consciência no que respeita a fazer
mantos, a recolher alimentos de esmola, a comer, a cuidar do kuti, o que fazer
nesta ou naquela situação. São todos conselhos muito práticos sobre a vida
diária de um bhikkhu. Um dia comum na vida de Bhikkhu Sumedho não é sobre
explodir em êxtase, mas sim levantar"-se, ir à casa de banho e vestir um roupão,
tomar banho, fazer isto ou aquilo; trata"-se apenas de estar consciente enquanto
se vive nesta forma e aprender a despertar para a forma como as coisas são, para
o Dhamma.

É por isso que sempre que contemplamos a cessação das coisas, não estamos à
procura do fim do universo, mas apenas do fim da expiração, da respiração, do
fim do dia, do fim do pensamento ou do fim do sentimento. Perceber isto,
significa que temos que prestar atenção ao fluxo da vida -- temos que realmente
perceber como ela é, em vez de esperar por algum tipo de experiência fantástica
de luz maravilhosa descendo sobre nós, electrocutando"-nos ou algo do género.

Contemplemos agora apenas a respiração normal do nosso corpo. Percebemos que, se
estivermos a inspirar, é fácil nos concentrarmos. Ao encher os pulmões, sentimos
uma sensação de crescimento, desenvolvimento e força. Quando dizemos que alguém
está emproado, provavelmente essa pessoa está inspirando. É difícil sentir"-se
emproado e altivo enquanto se expira. Ao expandirmos o nosso peito, teremos a
sensação de sermos alguém grande e poderoso. No entanto, quando comecei a
prestar atenção à expiração, a minha mente divagava; expirar não parecia tão
importante quanto inspirar -- estava apenas a fazer isso para poder começar a
próxima inspiração.

Agora reflictam nisto: é possível observar a respiração. E o que é que se pode
observar? O que é que observa e conhece a inspiração e a expiração? Não é a
respiração pois não? Também se pode observar o pânico que surge quando se quer
recuperar o fôlego e não se consegue, mas o observador, o que sabe, não é uma
emoção, não entra em pânico, não é uma expiração ou inspiração. Portanto, o
nosso refúgio no Buddha é ser esse saber, ser a testemunha e não a emoção,
respiração ou o corpo.

Com o som do silêncio\footnote{Ver capítulo `O Som do Silêncio'}, algumas
pessoas ouvem flutuações de som ou um fundo contínuo de som. Portanto podemos
contemplá"-lo, reparamos nisso -- será que podemos notá"-lo se pusermos os dedos
nos ouvidos? Conseguimos ouvi"-lo num local onde alguém está a utilizar uma
serra eléctrica? Ou quando se está a fazer exercícios ou num estado emocional
muito agitado? Estamos a usar este som do silêncio como algo para nos
lembrarmos, para onde nos dirigirmos e repararmos -- porque está sempre presente
aqui e agora. E existe aquilo que se apercebe dele.

Existe o desejo da mente para lhe chamar qualquer coisa, dar"-lhe um nome,
listá"-lo como espécie de realização ou de se lhe projectar algo. Reparem nisso,
na tendência para querer transformá"-lo em algo. Alguém disse que provavelmente
é apenas o som do sangue a circular nos ouvidos. Outra pessoa chamou"-lhe ``o
som cósmico'', ``a ponte para o Divino''. Isso soa melhor do que ``o sangue nos
teus ouvidos''. Pode ser o som do Cosmos ou pode ser que se tenha uma doença nos
ouvidos. Mas não tem de ser nada; é o que é, é ``simplesmente isso''. Seja o que
for, pode ser usado como reflexão, porque quando se está com isso, não há
sentido de eu, existe consciência, capacidade de reflectir.

Portanto, é mais como uma linha recta para a qual podemos ir, para evitar
prosseguirmos a cambalear. É algo que podemos usar para nos recompormos na vida
quotidiana, quando vestimos as vestes, quando lavamos os dentes, quando fechamos
uma porta, quando entramos na sala de meditação, quando nos sentamos. Grande
parte da vida quotidiana é apenas habitual, porque temos como objectivo o que
consideramos ser as coisas importantes da vida -- como a meditação. Para alguns,
caminhar desde onde vivem até à Sala de Meditação, pode ser uma experiência
totalmente descuidada -- apenas um hábito -- caminhando desconchavada e
ruidosamente! Depois, sentamo"-nos aqui durante uma hora a tentar estar
conscientes.

Desta forma, começamos a entrever uma maneira de sermos conscientes, de trazer
consciência às rotinas e experiências comuns da vida. Tenho um pequeno quadro no
meu quarto que gosto muito -- de um homem idoso com uma caneca de café na mão,
contemplando pela janela, um jardim inglês lá fora, com a chuva a cair. O título
do quadro é ``Esperando.'' Esta é a ideia que faço de mim próprio; um idoso com
a minha caneca de café, sentado ali à janela, esperando, esperando\ldots{}
observando a chuva ou o sol, seja o que for. Não vejo isto como uma imagem
depressiva, mas sim tranquila.

Esta vida limita"-se a esperar, não é? Esperamos o tempo tedo -- então
percebemos isso. Não estamos à espera de nada, mas podemos estar simplesmente à
espera. E então respondemos às coisas da vida, ao horário do dia, aos deveres, à
forma como as coisas se movem e mudam, à sociedade em que estamos. Essa resposta
não vem da força dos hábitos da cobiça, do ódio e da ilusão, mas é uma resposta
de sabedoria e consciência.

Agora, quantos de vocês sentem que têm uma missão na vida a cumprir? É algo que
têm de fazer, uma espécie de tarefa importante que vos foi atribuída por Deus,
pelo destino ou algo do género. As pessoas ficam frequentemente presas nessa
visão de serem alguém com uma missão. Quem consegue simplesmente aceitar as
coisas como elas são, de modo a que seja apenas o corpo que cresce, envelhece e
morre, que respira e está consciente? Podemos praticar, viver de acordo com os
preceitos morais, fazer o bem, responder às necessidades e experiências da vida
com atenção plena e sabedoria -- mas não há ninguém que tenha de \emph{fazer}
nada. Não há ninguém com uma missão, ninguém especial; não estamos a criar uma
pessoa, um santo, um avatar, um tulku, um messias ou um Maitreya. Mesmo que
penses: «Sou apenas um zé"-ninguém» -- mesmo ser um zé"-ninguém é ser alguém
nesta vida, não é? Podes ter tanto orgulho em ser um zé"-ninguém como em ser
alguém, e estar tão iludido ou apegado a ser um zé"-ninguém. Mas seja o que for
que acredites, que és um incómodo e um fardo para o mundo ou seja qual for a
forma como te vejas, então o conhecimento está lá para ver a cessação de tal
percepção.

As percepções surgem e cessam, não é verdade? «Sou alguém, uma pessoa importante
que tem uma missão na vida»: isso surge e cessa na mente. Repara no fim de ser
alguém importante ou de ser ninguém ou o que quer que seja -- tudo cessa, não é
verdade? Tudo o que surge, cessa, por isso não há apego à visão de ser alguém
com uma missão ou de ser ninguém. Há o fim de toda essa massa de sofrimento --
de ter de desenvolver algo, tornar"-se alguém, mudar algo, corrigir tudo,
livrar"-se de todas as suas impurezas ou salvar o mundo. Mesmo os melhores
ideais, os melhores pensamentos podem ser vistos como \emph{dhammas} que surgem
e cessam na mente.

Agora, pode pensar que esta é uma filosofia de vida árida, porque há muito mais
coração e sentimento em ser alguém que vai salvar todos os seres sencientes.
Pessoas com espírito de sacrifício, que têm missões, ajudam os outros e têm algo
importante a fazer, são uma inspiração. Mas quando percebe isso como
\emph{dhamma}, está a olhar para as limitações das inspirações e para a sua
cessação. Depois, há o \emph{dhamma} dessas aspirações e acções, em vez de
\emph{alguém} que tem de se tornar algo ou tem de fazer algo. Toda a ilusão é
abandonada e o que resta é a pureza da mente. Então, a resposta à experiência
vem da sabedoria e da pureza, em vez de vir da convicção pessoal e da missão com
o seu sentido de eu e do outro, e todas as complicações que advêm de todo esse
padrão de ilusão.

Consegues confiar nisso? Consegues confiar em simplesmente deixar tudo ir e
cessar, e não ser ninguém, não ter nenhuma missão, não ter de te tornares em
nada? Consegues realmente confiar nisso, ou achas isso assustador, árido ou
deprimente? Talvez queiras mesmo inspiração. `Diz"-me que está tudo bem; diz"-me
que me amas de verdade; que o que estou a fazer está certo e que o budismo não é
apenas uma religião egoísta onde se alcança a iluminação para o próprio bem;
diz"-me que o budismo está aqui para salvar todos os seres sencientes. É isso
que vais fazer, Venerável Sumedho? És realmente Mahayana ou Hinayana?'

O que estou a apontar é o que é a inspiração enquanto experiência. Idealismo:
não tentar descartá"-lo ou julgá"-lo de forma alguma, mas reflectir sobre ele,
para saber o que isso é na mente e com que facilidade podemos ser iludidos pelas
nossas próprias ideias e noções altivas. E ver o quão insensíveis, cruéis e
indelicados/as podemos ser devido ao apego que temos às noções sobre ser gentil
e sensível. É aqui que reside a verdadeira investigação do Dhamma.

Lembro"-me, na minha própria experiência, que sempre tive a visão de que era
alguém especial de alguma forma; costumava pensar: `Bem, \emph{devo} ser uma
pessoa especial. Há muito tempo, quando era criança, ficava fascinado pela Ásia
e, assim que pude, estudei chinês na universidade, por isso, certamente devo ter
sido a reencarnação de alguém que estava ligado ao Oriente.'

Mas considere isto como uma reflexão: não importa quantos sinais de ser especial
ou vidas passadas consiga recordar, ou vozes de Deus ou mensagens do Cosmos,
seja o que for -- não para negar isso ou dizer que essas coisas não são reais --
mas são impermanentes. São \emph{anicca}, \emph{dukkha, anatta}. Estamos a
reflectir sobre elas tal como realmente são -- o que surge cessa: uma mensagem
de Deus é algo que surge e cessa na tua mente, não é verdade? Deus não está
sempre a falar contigo continuamente, a menos que queiras considerar o silêncio
como a voz de Deus. Nesse caso, não diz realmente nada, pois não? Podemos
chamá"-lo de qualquer coisa: podemos chamá"-lo de voz de Deus, ou do divino, ou
do zumbido do cosmos, ou do sangue nos teus tímpanos. Mas seja o que for, pode
ser usado para a atenção plena e a reflexão -- é para isso que estou a apontar,
como usar estas coisas sem as transformar em algo.

Então, as missões que temos são respostas, não a experiências que temos nas
nossas vidas -- já não são pessoais, já não sou eu, Sumedho Bhikkhu, com uma
missão como se tivesse sido especialmente escolhido lá de cima, mais do que
qualquer um de vocês. Já não é isso. Toda essa maneira de pensar e perceber é
abandonada. E quer eu salve ou não o mundo e milhares de seres, ou ajude os
pobres nas favelas de Calcutá, ou ajude a curar todos os leprosos e faça todo o
tipo de boas obras -- não é a partir da ilusão de ser uma pessoa, é uma resposta
natural da sabedoria.

Nisto confio; isto é o que é \emph{saddha} -- é uma fé na palavra do Buddha.
\emph{Saddha}: é uma verdadeira confiança no Dhamma; em apenas esperar e não ser
ninguém e não me tornar em nada, mas ser capaz de apenas esperar e responder. E
se não houver muito a que responder, é apenas esperar -- chávena de café.
Observar a chuva, o pôr"-do-sol, envelhecer, testemunhar o processo de
envelhecimento, as idas e vindas no mosteiro -- as ordenações e as renúncias, as
inspirações e as depressões, os altos e os baixos, dentro da mente, fora no
mundo. E há a resposta porque, quando temos vigor para responder, é de uma forma
hábil e compassiva, que estamos muito dispostos e somos capazes de fazer.
Gostamos de ajudar as pessoas. Não me importaria de ir para uma colónia budista
de leprosos -- ficaria contente por isso -- ou de trabalhar nas favelas de
Calcutá ou em qualquer outro lugar, não teria objecções; esse tipo de coisas é
bastante apelativo para o meu sentido de nobreza!

Mas não é uma missão, não sou eu que tenho de fazer alguma coisa; é confiar no
Dhamma. Então, a resposta à vida é clara e benéfica, porque não vem de mim como
pessoa, e as ilusões, a inquietação, a compulsão, a obsessão da mente, cessam.
Largamos isso e acaba por cessar.

\photoPage{image-011-forgiveness.jpg}{Oferta formal a Tan Chao Khun Paññananda, Amaravati}{Foto de Jeff Pick}

